\section{Knowledge-representation patterns}\label{supp:kr}

This section gives formal patterns and connected graph examples for
the resource described in the main article. The Description Logic
notation \cite{dlhandbook} uses $\sqsubseteq$ for subclass inclusion,
$\sqcap$ for conjunction, $\exists R.C$ for at least one relation
$R$ to a member of class $C$, and $\forall R.C$ for restriction of
all $R$ fillers to $C$. The displayed patterns describe the intended
entity relationships; structural admission is checked separately.

Prepend Table~\ref{tab:supp_prefixes}'s declarations when parsing the
Turtle snippets. The unprefixed knowledge-base namespace is used by
the PowerSoil protocol. The listing validator checks examples against
the available module snapshot; full validation additionally requires
the corresponding taxonomy payload and pinned reference terms.

\begin{table}[h]
\centering
\caption{Prefix declarations used in all Turtle and SPARQL listings.}
\label{tab:supp_prefixes}
\small
\begin{tabular}{ll}
\toprule
\textbf{Prefix} & \textbf{Namespace} \\
\midrule
\texttt{rak:}      & \texttt{https://rubalkhali.science/kb/RAK\_} \\
\texttt{rak\_kb:}  & \texttt{https://rubalkhali.science/kb/} \\
\texttt{sio:}      & \texttt{http://semanticscience.org/resource/SIO\_} \\
\texttt{owl:}      & \texttt{http://www.w3.org/2002/07/owl\#} \\
\texttt{pato:}     & \texttt{http://purl.obolibrary.org/obo/PATO\_} \\
\texttt{uo:}       & \texttt{http://purl.obolibrary.org/obo/UO\_} \\
\texttt{envo:}     & \texttt{http://purl.obolibrary.org/obo/ENVO\_} \\
\texttt{ncbitaxon:}& \texttt{http://purl.obolibrary.org/obo/NCBITaxon\_} \\
\texttt{obo:}      & \texttt{http://purl.obolibrary.org/obo/} \\
\texttt{rdf:}      & \texttt{http://www.w3.org/1999/02/22-rdf-syntax-ns\#} \\
\texttt{rdfs:}     & \texttt{http://www.w3.org/2000/01/rdf-schema\#} \\
\texttt{xsd:}      & \texttt{http://www.w3.org/2001/XMLSchema\#} \\
\texttt{geo:}      & \texttt{http://www.opengis.net/ont/geosparql\#} \\
\texttt{dcterms:}  & \texttt{http://purl.org/dc/terms/} \\
\bottomrule
\end{tabular}
\end{table}

\subsection{Laboratory-control worked example}

The control module represents the whole-cell D6300 product, its
composition specification and evidence-bearing expected-taxon
assertions as separate entities. The specification links directly to
each supported taxon; a reified assertion additionally preserves the
subject, relation, taxon, status and source. Manufacturer genomic-DNA
mass fractions are retained in
\path{metadata/controls/control_composition.tsv}. Exact occurrence
identifiers and their source links are distributed in
\path{ontology/rubalkhali_controls.ttl}. This keeps long machine
identifiers in the executable resource while exposing the scientific
relationships here.

\subsection{Sampling pattern}

A sampling event is a process that derives a sample from a feature of
interest (the site).

\begin{align}
	\label{eq:sampling}
	\text{SamplingProcess} & \sqsubseteq \text{Process} \\
	\text{Sample} & \sqsubseteq \text{MaterialEntity} \\
    \text{SampleCollection} & \sqsubseteq \text{Collection} \sqcap
        \forall \text{hasMember}.\text{Sample} \\
	\text{SamplingProcess} & \sqsubseteq \exists
                                 \text{hasTarget}.\text{Site} \sqcap
                                 \exists
                                 \text{hasOutput}.\text{SampleCollection} \\
	\text{Sample} & \sqsubseteq \exists \text{isDerivedFrom}.\text{Site}
\end{align}

Listing~\ref{lst:ttl_sampling} is generated from the Trip~1, Site~1
deep-soil record in the development graph.

\begin{lstlisting}[caption={Turtle representation of a sampling process.}, label={lst:ttl_sampling}, basicstyle=\ttfamily\scriptsize, frame=single]
rak:P100001 a rak:0000024 ; # Sampling Process
  sio:000291 rak:1000001 ; # hasTarget
  sio:000229 rak:7000001 . # hasOutput (Collection)

rak:7000001 a sio:001419 ; # Material sample collection
  sio:000059 rak:6000001 . # hasMember

rak:6000001 a rak:0000021 ; # Deep Soil Sample
  sio:000244 rak:1000001 . # isDerivedFrom
\end{lstlisting}

The instance \texttt{rak:P100001} of a sampling
process targets the site \texttt{rak:1000001} and yields a
collection \texttt{rak:7000001} as its output. The collection
explicitly contains \texttt{rak:6000001}, and this deep-soil sample
is asserted to be derived
from the same site. The redundancy between \texttt{hasTarget} on the
process and \texttt{isDerivedFrom} on the sample is intentional: the
former records the action, the latter records the resulting
provenance edge directly on the material entity, so queries can
traverse from sample to site without joining through the sampling
event.

\subsection{Experimental protocol pattern}

A protocol is a reusable information specification. A protocol
execution is a process specified by that protocol. Each process (e.g., DNA extraction, library preparation)
consumes inputs and produces outputs, explicitly linking agents and
devices. The pattern instantiates SIO's general
\emph{Process--hasInput/hasOutput--Object} idiom and adds a
specification link to a reusable protocol description.

\begin{align}
	\label{eq:protocol}
	\text{ProtocolExecution} & \sqsubseteq \text{Process} \\
	\text{Protocol} & \sqsubseteq \text{InformationContentEntity} \\
	\text{ProtocolExecution} & \sqsubseteq \exists
	    \text{isSpecifiedBy}.\text{Protocol} \\
	\text{ProtocolExecution} & \sqsubseteq \exists
	    \text{hasInput}.\text{Entity} \sqcap \exists
	    \text{hasOutput}.\text{Entity}
\end{align}

Listing~\ref{lst:ttl_protocol} demonstrates the formalization
of a DNA extraction process.

\begin{lstlisting}[caption={Turtle representation of a DNA extraction protocol.}, label={lst:ttl_protocol}, basicstyle=\ttfamily\scriptsize, frame=single]
rak:P200001 a sio:000994 ; # Scientific Protocol Execution
  sio:000339 rak_kb:DNA_ext_powersoil ; # isSpecifiedBy
  sio:000230 rak:6000001 ; # hasInput (Soil Sample)
  sio:000229 rak:6800001 . # hasOutput (DNA Extract)

rak:6800001 a rak:0000040 ; # DNA Extract
  sio:000008 rak:5900001 . # hasAttribute (Concentration Quality)
\end{lstlisting}

The protocol-execution instance \texttt{rak:P200001}
is specified by the named extraction protocol
\texttt{rak\_kb:DNA\_ext\_powersoil}, takes the soil sample
\texttt{rak:6000001} as input, and produces the DNA extract
\texttt{rak:6800001} as output; the resulting extract carries a
concentration quality (\texttt{rak:5900001}, class
\texttt{rak:0000043}) which is itself
quantified through a separate measurement process following the
measurement pattern in the main article's Knowledge representation subsection. The measurement value that
quantifies it is the separate individual \texttt{rak:4900001}
(class \texttt{rak:0000044}), which points to the quality with
\texttt{sio:000215} and back from it with \texttt{sio:000216}.
Separate value and quality individuals preserve the reciprocal
measurement links.

\subsection{Taxonomic abundance pattern}

Taxonomic profiles contain read counts and relative read abundances
aggregated by rank. The validated mapping assigns NCBI or project
identifiers to each rank-truncated classifier lineage.
Listing~\ref{lst:ttl_tax} gives one such observation.

For instance, the sequence analysis for a specific profile yields an
absolute sequence count associated with an audited taxonomic class. We
construct an abundance quality that is an attribute of both that taxon
and the source sequence dataset, binding the count back to the
biological workflow. Formally, the pattern instantiates SIO's
\emph{Collection--hasMember--Thing} and reciprocal
\emph{MeasurementValue--isMeasurementValueOf--Quality} idioms:

\begin{small}
\begin{align}
	\label{eq:taxabundance}
	\text{BioinformaticWorkflow} & \sqsubseteq \text{Process} \\
	\text{TaxonomicDataset} & \sqsubseteq \text{Collection} \\
	\text{BioinformaticWorkflow} & \sqsubseteq \exists
	    \text{hasInput}.\text{SequenceDataset} \sqcap \exists
	    \text{hasOutput}.\text{TaxonomicDataset} \\
	\text{TaxonomicDataset} & \sqsubseteq \forall
	    \text{hasMember}.\text{AbundanceObservation} \\
	\text{AbundanceObservation} & \sqsubseteq \text{Quantity} \sqcap
	    \exists \text{isMeasurementValueOf}.\text{AbundanceQuality} \\
	\text{AbundanceQuality} & \sqsubseteq \exists
	    \text{isAttributeOf}.\text{Taxon} \\
	\text{AbundanceQuality} & \sqsubseteq \exists
	    \text{hasMeasurementValue}.\text{AbundanceObservation}
\end{align}
\end{small}

\begin{lstlisting}[caption={Turtle representation of taxonomic abundance formalized from a feature table.}, label={lst:ttl_tax}, basicstyle=\ttfamily\scriptsize, frame=single]
rak:P290001 a rak:0000071 ; # Bioinformatic Workflow
  sio:000339 rak:L000012 ; # 16S amplicon processing protocol
  sio:000230 rak:7791009 ; # hasInput (FASTQ Dataset, ERR16783308)
  sio:000229 rak:7740001 . # hasOutput (Taxonomic Dataset)

rak:7740001 a rak:0000074 ; # taxon absolute abundance dataset
  sio:000059 rak:V00000001 . # hasMember (Observation Value)

rak:V00000001 a rak:0000076 ; # Absolute abundance value
  sio:000215 rak:Q00000001 ; # isMeasurementValueOf
  rak:2000025 "Domain: Bacteria" ; # classifier lineage
  rak:2000026 "853814.0"^^xsd:double . # absolute read count

rak:Q00000001 a rak:0000078, pato:0000070 ; # Absolute Abundance Quality
  sio:000011 ncbitaxon:2, rak:7791009 ; # taxon and FASTQ bearers
  sio:000216 rak:V00000001 . # hasMeasurementValue
\end{lstlisting}

The bioinformatic-workflow instance
\texttt{rak:P290001} consumes the FASTQ dataset of run
\texttt{ERR16783308} (\texttt{rak:7791009}, profile \texttt{V10Dr1})
and produces the absolute taxonomic dataset
\texttt{rak:7740001}; that dataset has as one of its members the
measurement value \texttt{rak:V00000001}, which records an absolute
read count of \texttt{853{,}814} for the domain-rank record
\texttt{Domain: Bacteria}. Its reciprocal quality
\texttt{rak:Q00000001} is an attribute of both \emph{Bacteria}
(\texttt{ncbitaxon:2}) and the source FASTQ dataset. Each of the
fourteen datasets produced by one workflow instance holds a single
rank at one of the two abundance bases, so a rank-restricted query
never has to filter observations by lineage depth. Each taxon--rank entry adds a value and a quality while referencing
shared workflow, dataset and taxon objects. Contextual
lineages use the same pattern with their audited project IRI.

\subsection{XRF geochemistry pattern}

We model every XRF session as a process with multiple
concentration-value outputs and a workflow-specific protocol. A
laboratory session has an archived soil specimen as input. The field
instrument log identifies a test and site, and the graph links each
session to that site. Field-XRF and laboratory-XRF records remain
distinguishable even when they report the same analyte.
Analyte cross-references to ChEBI \cite{chebi} and PubChem
\cite{pubchem} are maintained in a single mapping ledger. The
identifier audit resolves every non-null identifier against pinned
ChEBI v247 and a dated PubChem snapshot, checking identifier
existence, label, formula, charge and atom-versus-oxide entity type.
Of 93 instrument channels, 79 have a supported ChEBI cross-reference
and 91 have a formula-matching PubChem cross-reference. Missing external identifiers remain explicitly null. The \texttt{rdfs:seeAlso} annotations identify chemical references
for reporting channels. The P$_2$O$_5$ channel reports oxide-equivalent
magnitudes; its ChEBI cross-reference P$_4$O$_{10}$ shares the
empirical ratio. Light Elements is an instrument-defined aggregate
channel with null chemical cross-references.

The shared output pattern specializes the measurement pattern
(main article, Knowledge representation); only the laboratory specialization
requires a soil-sample input:

\begin{align}
\label{eq:xrf}
\text{XRFProcess} & \sqsubseteq \text{MeasurementProcess}
  \sqcap \exists \text{hasOutput}.\text{Concentration} \\
\text{LabXRFProcess} & \sqsubseteq
  \text{XRFProcess} \sqcap \exists
  \text{hasInput}.\text{SoilSample} \\
\text{FieldXRFProcess} & \sqsubseteq
  \text{XRFProcess} \sqcap \exists
  \text{hasTarget}.\text{SamplingSite} \\
\text{Concentration} & \sqsubseteq \text{Quantity} \sqcap
  \exists \text{isMeasurementValueOf}.\text{AnalyteQuality}
\end{align}

Listing~\ref{lst:ttl_xrf} is generated from the aluminium output of
laboratory-XRF process \texttt{rak:P300072} in the current module.

\begin{lstlisting}[caption={Turtle representation of an XRF geochemistry measurement.}, label={lst:ttl_xrf}, basicstyle=\ttfamily\scriptsize, frame=single]
rak:P300072 a rak:0000025 ; # XRF Measuring Process
  sio:000230 rak:6001732 ; # hasInput (sample V1Sr1)
  sio:000339 rak:L_Surface_Lab ; # laboratory protocol
  sio:000229 rak:4302121 . # hasOutput

rak:4302121 a rak:0000507 ; # Al concentration value
  sio:000215 rak:5302121 ; # isMeasurementValueOf
  rak:2000012 "0.78"^^xsd:double . # Concentration

rak:5302121 a rak:0000107 ; # Al concentration quality
  sio:000011 rak:6001732 ; # isAttributeOf
  sio:000216 rak:4302121 . # hasMeasurementValue
\end{lstlisting}

This laboratory-XRF process consumes archived specimen
\texttt{rak:6001732} (sample \texttt{V1Sr1}) and produces an aluminium
concentration value.
The value points to its aluminium quality using
\texttt{sio:000215}; the quality points back using
\texttt{sio:000216}. A dedicated scan of every
XRF source artifact (the processed laboratory table, the field log
and all raw instrument exports)
(\path{scripts/xrf/audit_xrf_unit_evidence.py}) finds no unit
declaration attached to any concentration column; the only unit token
in the exports is \texttt{PPM}, and it belongs to the
lower-limit-of-detection column. The generated module therefore carries
no unit assertion for XRF concentrations. These values are
instrument-reported magnitudes with undocumented concentration units.
Field sessions use \texttt{sio:000291} to identify their site and
retain the instrument test identifier; laboratory sessions identify
their specimen input. Cross-workflow calibration requires documented
units, paired aliquots and reference materials.

\subsection{Sequencing workflow pattern}

Sequencing data are produced by a chain of laboratory and
computational processes that the knowledge graph represents as a
sequence of \emph{Process--hasInput/hasOutput--Object} steps: a DNA
extraction process consumes a soil sample and produces a DNA extract;
a library preparation process consumes the DNA extract and produces
an amplicon library; a sequencing process consumes the amplicon
library and produces a FASTQ dataset whose members are sequence reads.
Each process is specified by a protocol and may carry agents and
device participants, reusing the experimental protocol pattern in Equation~\ref{eq:protocol}.

\begin{small}
\begin{align}
	\label{eq:sequencing}
	\text{DNAExtractionProcess} & \sqsubseteq \text{ProtocolExecution} \\
	\text{DNAExtractionProcess} & \sqsubseteq \exists
	    \text{hasInput}.\text{SoilSample} \sqcap \exists
	    \text{hasOutput}.\text{DNAExtract} \\
	\text{LibraryPreparationProcess} & \sqsubseteq \text{ProtocolExecution} \\
	\text{LibraryPreparationProcess} & \sqsubseteq \exists
	    \text{hasInput}.\text{DNAExtract} \sqcap \exists
	    \text{hasOutput}.\text{AmpliconLibrary} \\
	\text{SequencingProcess} & \sqsubseteq \text{ProtocolExecution} \\
	\text{SequencingProcess} & \sqsubseteq \exists
	    \text{hasInput}.\text{AmpliconLibrary} \sqcap \exists
	    \text{hasOutput}.\text{FASTQDataset} \\
	\text{FASTQDataset} & \sqsubseteq \text{Dataset} \\
	\text{FASTQDataset} & \sqsubseteq \forall
	    \text{hasMember}.\text{SequenceRead}
\end{align}
\end{small}

Listing~\ref{lst:ttl_sequencing} instantiates the complete sequencing chain.

\begin{lstlisting}[caption={Turtle representation of the sequencing workflow from soil sample T1Dr1 to ENA run ERR16061083.}, label={lst:ttl_sequencing}, basicstyle=\ttfamily\scriptsize, frame=single]
rak:P200923 a sio:000994 ; # DNA Extraction Process
  sio:000339 rak_kb:DNA_ext_powersoil ; # isSpecifiedBy
  sio:000230 rak:6000577 ; # hasInput (T1Dr1 Soil Sample)
  sio:000229 rak:6800540 . # hasOutput (DNA Extract)

rak:P250001 a rak:0000065 ; # Library Preparation Process
  sio:000339 rak:L000010 ; # Illumina 16S Prep Guide
  sio:000230 rak:6800540 ; # hasInput (DNA Extract)
  sio:000229 rak:6690001 . # hasOutput (Amplicon Library)

rak:P280001 a rak:0000066 ; # Sequencing Process
  sio:000230 rak:6690001 ; # hasInput (Amplicon Library)
  sio:000229 rak:7790001 . # hasOutput (FASTQ Dataset)

rak:7790001 a rak:0000063 ; # FASTQ Dataset
  rdfs:seeAlso <https://identifiers.org/insdc.run/ERR16061083> .
\end{lstlisting}

The DNA extraction process
\texttt{rak:P200923} consumes the soil sample
\texttt{rak:6000577} and produces the DNA extract
\texttt{rak:6800540}; that extract is the input of the library
preparation \texttt{rak:P250001}, which yields the amplicon
library \texttt{rak:6690001}; finally the sequencing process
\texttt{rak:P280001} consumes the library and produces the FASTQ
dataset \texttt{rak:7790001}, which is cross-referenced to its
ENA accession via \texttt{rdfs:seeAlso}. Each link in the chain is
the same \texttt{hasInput}/\texttt{hasOutput} pattern, so a single
SPARQL traversal can recover the recorded provenance from any node in this
chain.

\subsection{Sequencing quality control pattern}

Sequencing quality control is represented as an additional
computational process that takes a FASTQ dataset as input and
produces multiple measurement values, each instantiating the
measurement pattern in the main article's Knowledge representation subsection.
Metrics (sequence count, GC content, duplicate rate, read length)
are split into forward (R1) and reverse (R2) subclasses so that the
provenance of paired-end runs is preserved.

\begin{small}
\begin{align}
	\label{eq:qc}
	\text{SequencingQCProcess} & \sqsubseteq \text{Process} \\
	\text{SequencingQCProcess} & \sqsubseteq \exists
	    \text{hasInput}.\text{FASTQDataset} \sqcap \exists
	    \text{hasOutput}.\text{QCMeasurementValue} \\
	\text{QCMeasurementValue} & \sqsubseteq \text{Quantity} \sqcap
	    \exists \text{isMeasurementValueOf}.\text{QCQuality} \\
	\text{QCQuality} & \sqsubseteq
	    \begin{aligned}[t]
	    &\exists \text{isAttributeOf}.{}\\[-0.25ex]
	    &\text{FASTQDataset}
	    \end{aligned}
\end{align}
\end{small}

Listing~\ref{lst:ttl_qc} shows the corresponding quality-control pattern.

\begin{lstlisting}[caption={Turtle representation of forward-read GC content for ENA run ERR16062319.}, label={lst:ttl_qc}, basicstyle=\ttfamily\scriptsize, frame=single]
rak:P350001 a rak:0000200 ; # Sequencing QC Process
  sio:000230 rak:7790263 ; # hasInput (FASTQ Dataset)
  sio:000229 rak:4450001 . # hasOutput (QC Value)

rak:4450001 a rak:0000217 ; # Forward GC Content Value
  sio:000215 rak:5550001 ; # isMeasurementValueOf
  sio:000221 uo:0000187 ; # Percent
  rak:2000074 "58.0"^^xsd:double . # has GC content

rak:5550001 a rak:0000205 ; # Forward GC Content Quality
  sio:000011 rak:7790263 ; # isAttributeOf (FASTQ Dataset)
  sio:000216 rak:4450001 . # hasMeasurementValue
\end{lstlisting}

The QC process \texttt{rak:P350001} takes the FASTQ
dataset \texttt{rak:7790263} as input and produces the
measurement value \texttt{rak:4450001}, which records a GC
content of \texttt{58.0}~percent and is the measurement value of a
forward-read GC-content quality \texttt{rak:5550001} that is in turn an
attribute of the same FASTQ dataset. Splitting the quality class into
forward and reverse subclasses preserves the read direction, so that
paired-end QC metrics can be queried independently per orientation.
